\section{Điều LIME chưa giải quyết}
\label{sec:ch03-dieu-lime-chua-giai-quyet}

Các tham số trên cho thấy LIME không trả lời xong câu hỏi về độ trung thực. Nó
cho một \tn{mô hình thay thế cục bộ}{local surrogate} có thể kiểm tra trên lân cận do người dùng dựng, nhưng
vẫn cần hỏi lân cận ấy có đại diện cho các thay đổi có nghĩa trong dữ liệu hay
không. Một đường khớp tốt trên các câu văn đã bị bỏ từ không tự chứng minh rằng
mô hình dựa vào cùng những từ đó trên văn bản tự nhiên.

Hơn nữa, tính thưa của $g$ là một ràng buộc về khả năng đọc. Nó không phải là bằng
chứng rằng các thành phần bị bỏ đi không có tác động. Khi $f$ có tương tác
phi tuyến ngay trong lân cận, một mô hình tuyến tính thưa có thể che các tương
tác ấy để đổi lấy một danh sách trọng số ngắn. Bài báo gốc cũng ghi nhận rằng
lớp mô hình tuyến tính có thể không đủ để giải thích một hành vi phi tuyến cục
bộ; khi đó chính bài báo đề nghị ước lượng điều nó gọi là faithfulness của lời
giải thích, tức độ khớp của $g$ với $f$ trên tập mẫu đã nhiễu, rồi trình con
số ấy cho người dùng~\autocite{p09lime}. Chữ faithfulness ở đó hẹp hơn
định nghĩa~\ref{def:ch01-do-trung-thuc}: nó đo $g$ bám $f$ đến đâu trên lân
cận đã dựng, không đo lân cận ấy có phản ánh cơ chế của $f$ hay không.

Do đó, LIME là một ví dụ tốt để phân biệt lời giải thích hữu ích với lời giải
thích đã được xác nhận. Nó đặt ra quy trình can thiệp rõ hơn việc chỉ tô đậm
một từ hay một vùng ảnh, nhưng bản thân quy trình vẫn chứa các giả định cần
được kiểm tra. Chương~\ref{ch:shap} sẽ đổi từ mô hình thay thế cục bộ sang
\tn{giá trị Shapley}{Shapley value}; phép đổi ấy không xóa câu hỏi về
\tn{phân phối nền}{background distribution} và về điều một
đóng góp thực sự đại diện.

\index{LIME|)}%
